% ============================================================================
% SECCIÓN 2.5: INVESTIGACIONES PREVIAS RELACIONADAS
% Generado por: Poseidón 🔱 (12 Nov 2025, 21:30 hrs)
% Narrativa: 700 palabras agrupadas por metodología
% ============================================================================

\section{Investigaciones Previas Relacionadas}
\label{sec:investigaciones_previas}

% ----------------------------------------------------------------------------
% INTRODUCCIÓN: Panorama del campo
% ----------------------------------------------------------------------------

La clasificación del comportamiento sedentario y la actividad física mediante dispositivos wearables ha experimentado un crecimiento exponencial en la última década, impulsado por la disponibilidad de sensores de consumo de bajo costo y el avance de las técnicas de inteligencia artificial. Los estudios recientes abordan este desafío desde múltiples perspectivas metodológicas, cada una con ventajas y limitaciones específicas en términos de precisión, interpretabilidad y generalización.

% ----------------------------------------------------------------------------
% SUBSECCIÓN: Enfoques de Deep Learning
% ----------------------------------------------------------------------------

\subsection{Enfoques de Deep Learning y Redes Neuronales}
\label{subsec:deep_learning}

Los enfoques de aprendizaje profundo han demostrado capacidades sobresalientes en el reconocimiento de patrones complejos en datos de sensores wearables. \cite{Khan2024Wearable} desarrollaron un sistema de reconocimiento de actividad física utilizando redes neuronales convolucionales (CNNs) y redes LSTM para procesar datos de sensores inerciales, logrando alta precisión en la clasificación de movimientos. Sin embargo, estos modelos enfrentan dos limitaciones críticas: requieren conjuntos de datos masivos para entrenamiento y operan como ``cajas negras'', dificultando la interpretación clínica de sus decisiones \cite{Bienefeld2023XAI}.

Un hallazgo particularmente relevante es el reportado por \cite{Bassani2025DLHAR}, quienes compararon el rendimiento de algoritmos de deep learning utilizando una división 70/30 frente a validación Leave-One-Subject-Out (LOSO). Los resultados mostraron una caída significativa del rendimiento: accuracy de 95.7\% en la división 70/30 versus 90.3\% en LOSO, evidenciando que la generalización robusta a nuevos usuarios requiere protocolos de validación más estrictos. Esta observación es fundamental para evaluar la aplicabilidad real de sistemas de clasificación en escenarios clínicos.

% ----------------------------------------------------------------------------
% SUBSECCIÓN: Machine Learning Clásico
% ----------------------------------------------------------------------------

\subsection{Machine Learning Clásico y Validación Cruzada}
\label{subsec:ml_clasico}

Los algoritmos de aprendizaje automático clásico, como Random Forest y Support Vector Machines, han sido ampliamente utilizados para la clasificación de actividades debido a su balance entre precisión y eficiencia computacional. \cite{Fuller2021Predicting} aplicaron Random Forest para clasificar posturas (lying, sitting, walking, running) utilizando datos de Apple Watch y Fitbit, alcanzando una precisión superior al 90\% mediante validación Leave-One-Out Cross-Validation (LOOCV) en una cohorte de 33 participantes. Este estudio validó el modelo Bring-Your-Own-Device (BYOD), demostrando la factibilidad de utilizar dispositivos de consumo en investigación clínica.

La importancia de la validación cruzada estricta ha sido enfatizada por \cite{Rehman2024LOSO}, quienes compararon el rendimiento de Random Forest utilizando LOSO versus k-fold cross-validation. Sus resultados mostraron que LOSO produjo un accuracy del 76\%, significativamente inferior al 89\% obtenido con k-fold, atribuyendo esta diferencia al \textit{data leakage} (fuga de datos) presente en k-fold cuando se aplica a datos de sensores wearables. Esta evidencia justifica la adopción de LOUO (Leave-One-User-Out) como estándar metodológico para evaluar la generalización inter-usuario.

El estudio de \cite{Mathew2024LOSO} refuerza esta conclusión al analizar el uso funcional de la mano en niños con parálisis cerebral unilateral mediante acelerómetros de muñeca. Los autores reportaron un F1-Score de 0.896 utilizando modelos personalizados por sujeto, que se redujo drásticamente a 0.584 al aplicar validación LOSO. Este fenómeno de sobreajuste es particularmente pronunciado en cohortes pequeñas ($N=22$), reflejando la dificultad de capturar la variabilidad inter-individual con tamaños de muestra limitados.

% ----------------------------------------------------------------------------
% SUBSECCIÓN: Lógica Difusa y Sistemas Expertos
% ----------------------------------------------------------------------------

\subsection{Lógica Difusa y Sistemas de Inferencia Interpretables}
\label{subsec:fuzzy_logic}

La lógica difusa emerge como una alternativa metodológica que prioriza la interpretabilidad sobre la precisión bruta, característica especialmente valiosa en aplicaciones biomédicas donde la confianza clínica depende de la transparencia del proceso de decisión. \cite{MarashiHosseini2023Dietary} desarrollaron un sistema de decisión dietética para pacientes con múltiples condiciones crónicas utilizando un sistema de inferencia difusa tipo Mamdani con 1144 reglas lingüísticas, alcanzando una precisión del 97\% al compararse con las recomendaciones de nutricionistas expertos ($N=100$). Este estudio demuestra que la lógica difusa puede igualar la precisión de enfoques estadísticos complejos mientras mantiene una estructura de conocimiento explícita y modificable.

La interpretabilidad total de los sistemas difusos ha sido destacada por \cite{Capitoli2025FuzzyXAI}, quienes proponen un modelo de lógica difusa probabilística para diagnóstico clínico que permite a los médicos ``generar, controlar y verificar todo el procedimiento de diagnóstico''. Este enfoque posiciona a la lógica difusa no solo como un clasificador, sino como una herramienta de Inteligencia Artificial Explicable (XAI) que genera confianza clínica mediante la exposición completa de su razonamiento.

Un precedente metodológico particularmente relevante para la presente investigación es el trabajo de \cite{Goncalves2021}, quienes desarrollaron un sistema de reconocimiento de actividad humana combinando clustering K-Means ($k=2$) con un sistema de inferencia difusa Mamdani. Este es el \textbf{único precedente identificado} en la literatura de un enfoque híbrido que utiliza clustering no supervisado para establecer una ground truth operativa, seguido de un sistema difuso supervisado. Este diseño metodológico resuelve el desafío de la ausencia de etiquetas gold-standard en datos de vida libre, al tiempo que preserva la interpretabilidad del modelo final.

% ----------------------------------------------------------------------------
% SUBSECCIÓN: Métodos Estadísticos Avanzados
% ----------------------------------------------------------------------------

\subsection{Métodos Estadísticos y Consensos Metodológicos}
\label{subsec:metodos_estadisticos}

Más allá de los enfoques de machine learning, se han desarrollado métodos estadísticos avanzados específicamente diseñados para la detección de patrones temporales en datos de wearables. \cite{Salim2024STEPHEN} propusieron un Hidden Semi-Markov Model para la detección de tiempo sedentario y \textit{sedentary bouts} (períodos continuos de sedentarismo) utilizando wearables de muñeca de consumo. Este enfoque estadístico captura la estructura temporal de las transiciones entre estados de actividad, ofreciendo una alternativa a los modelos de machine learning cuando la dinámica temporal es crítica.

La estandarización metodológica del campo ha sido abordada por \cite{Migueles2022GRANADA}, quienes presentan el consenso GRANADA sobre enfoques analíticos para evaluar asociaciones con comportamientos físicos determinados por acelerómetro en estudios epidemiológicos. Este consenso internacional establece las mejores prácticas para el análisis de datos de acelerómetro en investigación de actividad física, sedentarismo y sueño, proporcionando un marco de referencia para garantizar la comparabilidad entre estudios.

% ----------------------------------------------------------------------------
% SUBSECCIÓN: Validación de Dispositivos Wearables
% ----------------------------------------------------------------------------

\subsection{Validación de Precisión en Dispositivos Wearables Comerciales}
\label{subsec:validacion_wearables}

La precisión de los sensores wearables de consumo es un factor determinante para la validez de cualquier sistema de clasificación construido sobre sus datos. \cite{OGrady2024AppleWatch} realizaron un estudio de validación del Apple Watch Series 9 y Ultra 2 para la medición de variabilidad de frecuencia cardíaca (HRV) y frecuencia cardíaca en reposo (RHR), comparando contra el Polar H10 como gold standard en una cohorte de $N=39$ participantes. Los resultados mostraron un error porcentual absoluto medio (MAPE) de 28.88\% para HRV y 5.91\% para RHR, indicando que mientras la RHR es medida con alta precisión, la HRV presenta mayor variabilidad instrumental. Esta limitación debe considerarse al diseñar modelos que utilicen HRV como variable de entrada.

Un hallazgo conceptualmente relevante es el reportado por \cite{Godkin2025Context}, quienes demostraron que la RHR medida durante comportamiento sedentario es significativamente diferente de la RHR durante el sueño, reflejando contextos autonómicos distintos. Esta observación eleva la clasificación de sedentarismo de un problema postural (sentado vs acostado) a un problema fisiológico (estado autonómico sedentario vs estado de sueño), justificando la inclusión de variables cardiovasculares como HRV en sistemas de clasificación de sedentarismo.

\cite{Lyons2024StandHour} deconstruyeron el algoritmo de ``Stand Hour'' del Apple Watch, identificando los umbrales específicos de energía activa y conteo de pasos que activan esta métrica comercial. Este trabajo es crucial para entender los algoritmos de caja negra presentes en dispositivos de consumo y proporciona un punto de comparación para la interpretabilidad de modelos abiertos basados en lógica difusa.

% ----------------------------------------------------------------------------
% SUBSECCIÓN: HRV y Cognición
% ----------------------------------------------------------------------------

\subsection{Relación entre HRV, Actividad Física y Salud}
\label{subsec:hrv_cognicion}

La variabilidad de frecuencia cardíaca (HRV) no solo es un indicador del estado autonómico, sino que ha sido asociada con múltiples resultados de salud. \cite{Marino2024ARIC} analizaron datos de 961 adultos mayores del estudio ARIC Neurocognitive, encontrando que mayor actividad física y mayor HRV se asocian con mejor función cognitiva, con una relación no lineal entre actividad física, HRV y salud cerebral. Este hallazgo respalda la hipótesis de que la HRV captura información fisiológica complementaria a las métricas de movimiento, justificando su inclusión en modelos de clasificación de comportamiento sedentario.

% ============================================================================
% FIN SECCIÓN 2.5
% ============================================================================

